\section*{Capítulo \ref{ch:b1:heat-budget} --- El balance térmico de la Tierra}

\begin{solution}{exo:b1:heat-budget:1}
$q = k\,\mathrm dT/\mathrm dz = 2.5\times\qty{0.022}{K.m^{-1}} =
\qty{0.055}{W.m^{-2}} = \qty{55}{mW.m^{-2}}$. A través de \qty{1}{km^2} $=
\qty{e6}{m^2}$: \qty{5.5e4}{W}, es decir \qty{55}{kW}.
\end{solution}

\begin{solution}{exo:b1:heat-budget:2}
Continentes $\num{2.073e14}\times0.0709 = \qty{1.470e13}{W} = \qty{14.70}{TW}$;
océanos $\num{3.028e14}\times0.1054 = \qty{31.92}{TW}$; total \qty{46.61}{TW}, y
el \omterm{def:b1:heat-budget:heat-flow}{flujo de calor} medio $\num{4.661e13}/\num{5.101e14} = \qty{91.4}{mW.m^{-2}}$ (los
valores publicados, 46.7 y 91.6, difieren solo por el redondeo de las dos medias). La
fracción de la entrada solar es $46.6/173000 = \num{2.7e-4}$, o
\qty{0.027}{\percent}.
\end{solution}

\begin{solution}{exo:b1:heat-budget:3}
Uranio: $0.99274\times\num{9.513e-5} + 0.00720\times\num{5.685e-4} =
\num{9.444e-5} + \num{4.09e-6} = \qty{9.85e-5}{W.kg^{-1}}$; el raro
\ce{^{235}U} aporta el \qty{4}{\percent} de ella. Potasio: la fracción másica de
\ce{^{40}K} es $\num{1.17e-4}\times39.96/39.10 = \num{1.196e-4}$, así que
$\num{1.196e-4}\times\num{2.847e-5} = \qty{3.40e-9}{W.kg^{-1}}$.
\end{solution}

\begin{solution}{exo:b1:heat-budget:4}
Conversión: $\qty{1}{\micro cal.cm^{-2}.s^{-1}} = \num{4.184e-6}\,\unit{J.s^{-1}}/\qty{e-4}{m^2}
= \qty{41.84}{mW.m^{-2}}$: Estados Unidos orientales \qty{33.5}{mW.m^{-2}}, Basin and
Range \qty{58.6}{mW.m^{-2}}, Sierra Nevada \qty{16.7}{mW.m^{-2}}. Con $A_0D$ en
\unit{mW.m^{-2}} cuando $A_0$ está en \unit{\micro W.m^{-3}} y $D$ en \unit{km}:
$33.5 + 3.0\times7.5 = \qty{56.0}{mW.m^{-2}}$, $58.6 + 3.0\times9.4 =
\qty{86.8}{mW.m^{-2}}$, $16.7 + 3.0\times10 = \qty{46.7}{mW.m^{-2}}$. El mismo
granito da \omterm{def:b1:heat-budget:heat-flow}{flujos de calor} casi el doble de altos en Basin and Range que en la
Sierra Nevada: la diferencia está en el \omterm{def:b1:layered-earth:layers}{manto} de abajo.
\end{solution}

\begin{solution}{exo:b1:heat-budget:5}
$T(10\,\text{km}) = 10 + \frac{0.060\times10^4}{2.5} - \frac{\num{2e-6}\times10^8}{2\times2.5}
= 10 + 240 - 40 = \qty{210}{\degreeCelsius}$. Debajo, $q_m = 60 - 2.0\times10 =
\qty{40}{mW.m^{-2}}$, un gradiente de \qty{16}{K.km^{-1}}, y
$T(35\,\text{km}) = 210 + 16\times25 = \qty{610}{\degreeCelsius}$.
\end{solution}

\begin{solution}{exo:b1:heat-budget:6}
Lutita: gradiente $12.0/0.300 = \qty{40}{K.km^{-1}}$, $q = 1.5\times0.040 =
\qty{60}{mW.m^{-2}}$. Granito: $10.0/0.500 = \qty{20}{K.km^{-1}}$,
$q = 3.0\times0.020 = \qty{60}{mW.m^{-2}}$. $R(300) = 300/1.5 =
\qty{200}{m^2.K.W^{-1}}$, $R(800) = 200 + 500/3.0 = \qty{366.7}{m^2.K.W^{-1}}$.
Los tres puntos $(0, 12.0)$, $(200, 24.0)$, $(366.7, 34.0)$ están alineados, con
pendiente $22.0/366.7 = \qty{0.0600}{W.m^{-2}}$: el \omterm{def:b1:heat-budget:heat-flow}{flujo de calor} es el mismo,
\qty{60}{mW.m^{-2}}, en ambas rocas; solo cambia el gradiente, en proporción
inversa a la conductividad.
\end{solution}

\begin{solution}{exo:b1:heat-budget:7}
$\tau = L^2/\kappa$: \qty{1}{m}, \qty{e6}{s} o \num{12}~días; \qty{10}{m},
\qty{e8}{s} o \num{3.2}~años; \qty{100}{km}, \qty{e16}{s} o
\qty{317}{Myr}; \qty{2900}{km}, \qty{8.4e18}{s} o \qty{267}{Gyr}. A
\qty{10}{m} el ciclo diario no se siente en absoluto, y el anual, más corto
que los tres años necesarios para atravesar \qty{10}{m}, solo muy débilmente: solo llegan
los cambios que duran años o más. El último tiempo es $267/4.54 = 59$ veces la
edad de la Tierra.
\end{solution}

\begin{solution}{exo:b1:heat-budget:8}
(a) $1646 - 0.54\times50 = \qty{1619}{K}$, es decir \qty{1346}{\degreeCelsius},
el \qty{1350}{\degreeCelsius} publicado, redondeado. (b) Gradiente medio
$(0.54 + 0.36)/2 = \qty{0.45}{K.km^{-1}}$, así que $1646 + 0.45\times360 =
\qty{1808}{K}$, \qty{9}{K} (\qty{0.5}{\percent}) por encima de \qty{1799}{K}: el
gradiente no disminuye linealmente, pero la estimación es cercana.
\end{solution}

\begin{solution}{exo:b1:heat-budget:9}
(a) $G^2 = T_0^2/(\pi\kappa t)$ da $t = T_0^2/(\pi\kappa G^2)$.
(b) $G = 0.02\times\tfrac59\,\unit{K}/\qty{0.3048}{m} = \qty{0.03645}{K.m^{-1}} =
\qty{36.45}{K.km^{-1}}$; $t = 3900^2/(\pi\times\num{1.2e-6}\times0.03645^2) =
\num{1.521e7}/\num{5.009e-9} = \qty{3.04e15}{s}$, es decir \qty{96}{Myr},
la cifra de Kelvin. (c) $t = 1300^2/(\pi\times\num{e-6}\times0.036^2) =
\qty{4.15e14}{s} = \qty{13}{Myr}$: una temperatura inicial más realista hace
que la Tierra conductora sea aún más joven.
\end{solution}

\begin{solution}{exo:b1:heat-budget:10}
(a) Tras la erosión, la nueva superficie está a la antigua profundidad $e$, donde la
producción de calor era $A_0' = A_0e^{-e/D}$ y el \omterm{def:b1:heat-budget:heat-flow}{flujo de calor} $q(e) = q_m +
A_0De^{-e/D} = q_m + A_0'D$; medida desde la nueva superficie, la producción
sigue disminuyendo como $A_0'e^{-z'/D}$. El plutón se desplaza hacia abajo a lo largo de la misma recta.
(b) La capa tiene ahora un espesor $h - e$ con la misma $A$, así que $q_0' = q_m + A(h - e)$:
sobre una recta de pendiente $h - e$ que pasa por $q_m$. (c) Solo el modelo exponencial: en el
modelo de capa uniforme, los plutones erosionados a profundidades distintas se situarían sobre rectas de
pendientes diferentes, y los puntos de una provincia se dispersarían. Que se observe una
sola recta favorece una producción que disminuye aproximadamente de forma exponencial con la
profundidad.
\end{solution}

\begin{solution}{exo:b1:heat-budget:11}
(a) $\tau_P = RL/(3\kappa)$: para \qty{50}{km}, $\num{6.4e6}\times\num{5e4}/\num{3e-6}
= \qty{1.07e17}{s} = \qty{3.4}{Gyr}$; para \qty{100}{km}, \qty{6.8}{Gyr}.
(b) $q = kT_i/L = 3\times1300/\num{5e4} = \qty{78}{mW.m^{-2}}$ al principio,
y $78\,e^{-2/3.38} = \qty{43}{mW.m^{-2}}$ tras \qty{2}{Gyr}. (c)
$3\times1300/\sqrt{\pi\times\num{e-6}\times\num{6.31e16}} = 3900/\num{4.45e5} =
\qty{8.8}{mW.m^{-2}}$: la Tierra bien agitada sigue perdiendo cinco veces más calor
tras dos mil millones de años. Todo su interior alimenta la superficie, mientras que en una
Tierra conductora solo lo hace una cáscara exterior.
\end{solution}

\begin{solution}{exo:b1:heat-budget:12}
(a) $\sqrt{\pi\kappa\times\qty{1}{Myr}} = \sqrt{\pi\times\num{e-6}\times\num{3.156e13}}
= \qty{9957}{m}$, así que $C = 3\times1336/9957 = \qty{0.403}{W.m^{-2}}$ para
$t$ en Myr, \qty{403}{mW.m^{-2}.Myr^{1/2}}: el \qty{82}{\percent} de los 490 ajustados. El
modelo necesita un $k\,\Delta T$ mayor (unos \qty{4900}{W.m^{-1}}): una conductividad
mayor del \omterm{def:b1:layered-earth:layers}{manto}, o un interior más caliente, como a lo largo de la adiabática.
(b) $490/\sqrt1 = 490$, $490/\sqrt{10} = 155$ y $490/\sqrt{100} =
\qty{49}{mW.m^{-2}}$. (c) $z = 2.326\sqrt{\num{e-6}\times\num{3.156e15}} =
2.326\times\qty{56180}{m} = \qty{131}{km}$.
\end{solution}

\begin{solution}{pb:b1:heat-budget:1}
\textbf{1.} Gradiente $= (T_{500} - T_{100})/\qty{0.4}{km}$: A \qty{15.25}{K.km^{-1}},
B \num{19.50}, C \num{19.50}, D \num{25.50}, E \qty{28.25}{K.km^{-1}}.

\textbf{2.} Perforar altera la temperatura de la roca (el fluido de perforación enfría o calienta las paredes), y la perturbación
tarda semanas o meses en desaparecer. Las decenas de metros superiores siguen el ciclo estacional y los cambios recientes de la
temperatura en superficie.

\textbf{3.} $q_0 = k\times$gradiente: A $2.9\times15.25 = \qty{44.225}{mW.m^{-2}}$,
B \qty{50.70}{mW.m^{-2}}, C \qty{60.45}{mW.m^{-2}}, D \qty{68.85}{mW.m^{-2}},
E \qty{84.75}{mW.m^{-2}}.

\textbf{4.} De A a D están por debajo de la media continental, y E por encima en
\qty{13.9}{mW.m^{-2}}, o \qty{20}{\percent}.

\textbf{5.} El \omterm{def:b1:heat-budget:heat-flow}{flujo de calor} es la conductividad por el gradiente: C conduce mejor (3.1
frente a 2.6), así que el mismo gradiente transporta un \qty{19}{\percent} más de calor.

\textbf{6.} Las isotermas siguen el relieve solo en parte: bajo una cima se elevan menos que la superficie, así que están más
separadas por debajo, y el gradiente medido disminuye (aumenta bajo un valle). Hace falta una corrección por relieve antes de
comparar los plutones.

\textbf{7.} $\rho C_U H_U$: A $2650\times\num{2.0e-6}\times\num{9.85e-5} =
\qty{0.522}{\micro W.m^{-3}}$, B 0.914, C 1.305, D 1.827, E 2.610.

\textbf{8.} Torio: A 0.558, B 0.976, C 1.255, D 1.742, E 2.439.

\textbf{9.} Potasio: A $2650\times0.030\times\num{3.40e-9} = 0.270$, B 0.315,
C 0.360, D 0.378, E 0.405.

\textbf{10.} $A_0$ = 1.35, 2.20, 2.92, 3.95 y \qty{5.46}{\micro W.m^{-3}}.
El potasio aporta lo menos, desde el \qty{20}{\percent} del total en A hasta
\qty{7}{\percent} en E; el uranio y el torio se reparten el resto casi por igual.

\textbf{11.} Sí: los cinco puntos suben de forma regular, a unos
\qty{1.5}{mW.m^{-2}} de una recta.

\textbf{12.} $\bar A_0 = 15.88/5 = 3.176$, $\bar q_0 = 308.975/5 = 61.795$;
$\sum(A_0 - \bar A_0)^2 = 10.1681$; $\sum(A_0 - \bar A_0)(q_0 - \bar q_0) =
101.146$.

\textbf{13.} $D = 101.146/10.1681 = 9.95$ (en \unit{mW.m^{-2}} por
\unit{\micro W.m^{-3}}, es decir en km): $D = \qty{9.95}{km}$, y $q_m = 61.795 -
9.95\times3.176 = \qty{30.20}{mW.m^{-2}}$.

\textbf{14.} La pendiente, de unos \qty{10}{km}, es la de las tres provincias; el \omterm{def:b1:heat-budget:heat-flow}{flujo de calor} reducido, \qty{30}{mW.m^{-2}},
está cerca de los \qty{33}{mW.m^{-2}} de los Estados Unidos orientales: una provincia tranquila con un \omterm{def:b1:layered-earth:layers}{manto} frío debajo, a
diferencia de Basin and Range.

\textbf{15.} $q_m/k = \qty{12}{K.km^{-1}}$ y $A_0D^2/k = \num{5.46e-6}\times\num{e8}/2.5
= \qty{218.4}{K}$: $T(z) = 10 + 12z + 218.4\,(1 - e^{-z/10})$, $z$ en km, $T$ en
\unit{\degreeCelsius}.

\textbf{16.} $q(z) = 30 + 54.6\,e^{-z/10}$: \qty{50.1}{mW.m^{-2}} a \qty{10}{km},
\num{37.4} a \qty{20}{km} y \qty{31.7}{mW.m^{-2}} a \qty{35}{km}.

\textbf{17.} $30/84.6 = \qty{35}{\percent}$: casi dos tercios del \omterm{def:b1:heat-budget:heat-flow}{flujo de calor} sobre E se producen en su capa radiogénica.

\textbf{18.} $T(10) = 10 + 120 + 218.4\times0.632 = \qty{268}{\degreeCelsius}$.

\textbf{19.} $T(20) = 10 + 240 + 218.4\times0.865 = \qty{439}{\degreeCelsius}$.

\textbf{20.} Bajo A, $A_0D^2/k = 1.35\times100/2.5 = \qty{54.0}{K}$:
$T(35) = 10 + 420 + 54.0\times0.970 = \qty{482}{\degreeCelsius}$.

\textbf{21.} Por $A_0D^2/k$: \qty{54}{K} bajo A y \qty{218}{K} bajo E. El \omterm{def:b1:heat-budget:heat-flow}{flujo de calor} profundo es el mismo; la \omterm{def:b1:layered-earth:layers}{corteza} de E es
cuatro veces más rica en fuentes de calor, y el calor que producen debe conducirse hacia arriba por la misma roca, lo que exige un
gradiente más empinado en todo el recorrido.

\textbf{22.} $1300 = 10 + 218.4 + 12z$ da $z = \qty{89.3}{km}$ bajo E;
$1300 = 10 + 54.0 + 12z$ da \qty{103.0}{km} bajo A.

\textbf{23.} Ambas se encuentran dentro del rango de la \omterm{def:b1:layered-earth:lithosphere}{litosfera} continental, desde cien kilómetros o menos bajo regiones jóvenes
hasta unos \qty{250}{km} bajo cratones, hacia su lado delgado: con un \omterm{def:b1:heat-budget:heat-flow}{flujo de calor} reducido de \qty{30}{mW.m^{-2}} y esta
conductividad, el \omterm{def:b1:heat-budget:geotherm}{geotermo} conductivo alcanza temperaturas del \omterm{def:b1:layered-earth:layers}{manto} a unos \qty{100}{km}; bajo el plutón más radiactivo,
\qty{14}{km} antes.

\textbf{24.} $T(35) = 10 + 420 + 218.4\times(1 - e^{-3.5}) = 430 + 211.8$: el
\omterm{def:b1:layered-earth:discontinuity}{Moho} bajo el plutón E está a
\textbf{\qty[text-series-to-math=true]{642}{\degreeCelsius}}.
\end{solution}
